\documentclass[12pt,a4paper]{article}

% -------------------- Packages --------------------
\usepackage[a4paper,margin=1in]{geometry}
\usepackage{setspace}
\usepackage{titlesec}
\usepackage{enumitem}
\usepackage{array}
\usepackage{hyperref}
\usepackage{fancyhdr}

% -------------------- Page Formatting --------------------
\onehalfspacing
\setlength{\parindent}{0pt}
\setlength{\parskip}{6pt}

\titleformat{\section}
{\large\bfseries}
{\thesection.}
{0.5em}
{}

\titleformat{\subsection}
{\normalsize\bfseries}
{\thesubsection}
{0.5em}
{}

\setlist[itemize]{
    leftmargin=1.5em,
    itemsep=3pt
}

% -------------------- Header / Footer --------------------
\pagestyle{fancy}
\fancyhf{}
\lhead{\textbf{Journal 3}}
\rhead{\textbf{Smart City Civic Complaint System}}
\cfoot{\thepage}

% =========================================================
\begin{document}
% =========================================================

\begin{center}
    {\Large \textbf{JOURNAL 3}}\\[8pt]
    {\large \textbf{Frontend Implementation}}\\[6pt]
    {\normalsize \textbf{Smart City Civic Complaint and Management System}}
\end{center}

\vspace{12pt}

\begin{tabular}{|p{0.28\textwidth}|p{0.62\textwidth}|}
\hline
\textbf{Project} &
Smart City Civic Complaint and Management System \\ \hline

\textbf{Development Stage} &
Frontend Prototype Implementation \\ \hline

\textbf{Technology Used} &
React, Vite, React Router, JavaScript, CSS \\ \hline

\textbf{State Management} &
React Context API \\ \hline


\textbf{Student} &
Hittin \\ \hline
\end{tabular}

% =========================================================
\section{Objective}
% =========================================================

The work primarily focused on developing the Field Worker interface and the
Administrator interface. The Field Worker interface was designed to support
task acceptance, work progress tracking, evidence submission, and task
completion. The Administrator interface was developed to provide an overview
of complaints and their current processing status.

% =========================================================
\section{Development Environment and Technologies}
% =========================================================

The frontend components were developed using the following technologies:

\begin{itemize}
    \item \textbf{React} -- used for creating reusable and interactive
    frontend components.
    
    \item \textbf{Vite} -- used as the frontend development environment and
    build tool.
    
    \item \textbf{JavaScript} -- used for implementing component logic and
    user interactions.
    
    \item \textbf{React Router} -- used for navigation between dashboard and
    task-related pages.
    
    \item \textbf{CSS} -- used for designing dashboards, cards, buttons,
    progress indicators, forms, and status displays.
    
    \item \textbf{React Context API} -- used for accessing and updating
    shared complaint information between frontend components.
\end{itemize}

% =========================================================
\section{Field Worker Dashboard}
% =========================================================

The Field Worker Dashboard was developed to provide field workers with a
centralized interface for viewing their assigned civic complaint tasks.

The dashboard provides information about assigned work and the current
progress of tasks.

The main sections include:

\begin{itemize}
    \item Assigned Tasks
    \item In Progress Tasks
    \item Completed Tasks
    \item Complaint task information
    \item Task status
    \item Location information
    \item Task details navigation
\end{itemize}

Each task provides a mechanism to open its detailed task interface.

This allows the field worker to move from the dashboard to the individual
task workflow.

% =========================================================
\section{Worker Task Details}
% =========================================================

A separate \textit{Worker Task Details} page was developed for handling the
complete lifecycle of an assigned field task.

The task details interface displays important information about the civic
complaint and provides controls for updating the work status.

The implemented workflow is:

\begin{center}
\textbf{Assigned}
$\rightarrow$
\textbf{Accepted}
$\rightarrow$
\textbf{In Progress}
$\rightarrow$
\textbf{Completed}
\end{center}

Each stage represents a different point in the field-worker workflow.

% =========================================================
\section{Accept Task}
% =========================================================

An \textit{Accept Task} operation was implemented for field workers.

When a worker accepts an assigned task, the task status changes from
\textit{Assigned} to \textit{Accepted}.

This action represents that the field worker has received and accepted the
responsibility for performing the assigned work.

The interface dynamically reflects the updated status after the action.

% =========================================================
\section{Start Work}
% =========================================================

After accepting a task, the field worker can start the assigned work.

The \textit{Start Work} operation changes the task status to
\textit{In Progress}.

Initial progress is also set when work begins.

This provides a clear distinction between a task that has only been accepted
and a task on which actual field work has started.

% =========================================================
\section{Work Progress Tracking}
% =========================================================

A progress tracking section was implemented to allow the field worker to
update the percentage of work completed.

A progress slider was provided in the task interface.

The progress value represents the current completion level of the field
work.

For example:

\begin{itemize}
    \item 0\% -- Work has not started.
    \item 10\% -- Work has recently started.
    \item 50\% -- Approximately half of the work is completed.
    \item 100\% -- Work is completely finished.
\end{itemize}

The interface also contains a visual progress bar so that the current
completion percentage can be understood easily.

% =========================================================
\section{Work Notes}
% =========================================================

A work-notes section was implemented to allow the field worker to provide
additional information about the performed work.

The worker can enter notes describing activities carried out during the
resolution process.

Examples of information that can be recorded include:

\begin{itemize}
    \item Work performed.
    \item Issue identified at the location.
    \item Materials or actions used for resolution.
    \item Additional observations.
\end{itemize}

These notes provide supporting information before the task is marked as
completed.

% =========================================================
\section{Evidence Upload}
% =========================================================

An evidence-upload interface was implemented for the field worker.

The worker can select an evidence file after carrying out the assigned work.

The interface displays the selected file so that the worker can verify that
evidence has been attached.

Evidence is intended to support the completion of a civic complaint by
providing a record of the work performed.

At the current frontend stage, the evidence file is handled through the
frontend interface and is not permanently stored in a backend database.

% =========================================================
\section{Task Completion Validation}
% =========================================================

Validation was implemented before allowing a field worker to complete a
task.

The task cannot be marked as completed unless the required work information
has been provided.

The completion condition requires:

\begin{itemize}
    \item Work progress to reach 100\%.
    \item Work notes to be entered.
\end{itemize}

This prevents a task from being marked as completed prematurely.

The validation provides a basic frontend-level control for maintaining the
correct task workflow.

% =========================================================
\section{Complete Task}
% =========================================================

After the work progress reaches 100\% and the required work notes have been
provided, the field worker can select the \textit{Complete Task} operation.

The task status is then changed to:

\begin{center}
\textbf{Completed}
\end{center}

The completed-task interface provides a clear indication that the field
work has reached its final stage.

This status can subsequently be used by the complaint workflow to indicate
that the resolution is ready for further processing.

% =========================================================
\section{Worker Interface Styling}
% =========================================================

CSS styling was implemented for the Field Worker interface to improve
readability and usability.

The styling covers:

\begin{itemize}
    \item Worker dashboard layout.
    \item Task cards.
    \item Status indicators.
    \item Progress slider.
    \item Progress bar.
    \item Work-notes section.
    \item Evidence-upload section.
    \item File selection display.
    \item Task action buttons.
    \item Completed-task interface.
\end{itemize}

The interface was tested in the browser to verify that the different task
states and controls were displayed correctly.

% =========================================================
\section{Administrator Dashboard}
% =========================================================

An Administrator Dashboard was developed to provide an overall view of the
complaint management system.

The administrator can view complaint information and monitor the current
state of complaints through a centralized dashboard.

The dashboard displays complaint-related statistics and a complaint list.

% =========================================================
\section{Administrator Statistics}
% =========================================================

Dynamic complaint statistics were implemented in the Administrator
Dashboard.

The dashboard calculates the number of complaints belonging to different
status categories.

The implemented statistics include:

\begin{itemize}
    \item Total Complaints
    \item Pending Complaints
    \item In Progress Complaints
    \item Completed Complaints
    \item Resolved Complaints
    \item Reopened Complaints
\end{itemize}

These values are calculated from the complaint collection available through
the shared frontend state.

Therefore, the displayed values change when complaint information changes.

% =========================================================
\section{Administrator Complaint Management View}
% =========================================================

A complaint listing was implemented in the Administrator Dashboard.

For each complaint, the administrator can view important information such
as:

\begin{itemize}
    \item Complaint ID
    \item Complaint title
    \item Complaint category
    \item Location
    \item Submission date
    \item Current status
    \item Assigned worker
\end{itemize}

This provides the administrator with a consolidated view of complaints
being handled by the system.

% =========================================================
\section{Shared State Integration}
% =========================================================

The Administrator Dashboard was connected with the shared complaint state
implemented using React Context API.

The dashboard obtains the complaint collection using the
\texttt{useComplaints()} hook.

The shared state contains information such as:

\begin{itemize}
    \item Complaint status
    \item Assigned worker
    \item Work progress
    \item Work notes
    \item Evidence information
    \item Citizen verification status
\end{itemize}

This allows the administrator interface to display the current frontend
state of complaints rather than maintaining a completely separate complaint
dataset.

% =========================================================
\section{Frontend Navigation}
% =========================================================

The Field Worker and Administrator interfaces were connected to the
application routing structure.

The major routes associated with these interfaces include:

\begin{itemize}
    \item Worker Dashboard
    \item Worker Task Details
    \item Administrator Dashboard
\end{itemize}

React Router was used to navigate between these pages.

The worker can open an individual task from the dashboard, while the
administrator can access the centralized complaint overview.

% =========================================================
\section{Testing and Debugging}
% =========================================================

The implemented interfaces were tested through the React development
server.

The following operations were verified:

\begin{itemize}
    \item Opening the Field Worker Dashboard.
    \item Opening individual worker tasks.
    \item Accepting an assigned task.
    \item Starting a task.
    \item Updating work progress.
    \item Entering work notes.
    \item Selecting evidence files.
    \item Completing a task after satisfying validation conditions.
    \item Opening the Administrator Dashboard.
    \item Displaying complaint statistics.
    \item Displaying complaint information in the administrator view.
    \item Navigating between the implemented frontend routes.
\end{itemize}

CSS and component-level issues encountered during development were corrected
and the interfaces were verified in the browser.

% =========================================================
\section{Current Frontend Result}
% =========================================================

The Field Worker and Administrator frontend modules were successfully
implemented as part of the individual contribution.

The Field Worker module supports the complete prototype task progression:

\begin{center}
\textbf{Assigned}
$\rightarrow$
\textbf{Accepted}
$\rightarrow$
\textbf{In Progress}
$\rightarrow$
\textbf{Completed}
\end{center}

The Administrator module provides a centralized overview of complaints,
including dynamic complaint statistics and complaint status information.

These modules form the staff-side portion of the frontend prototype.

% =========================================================
\section{Current Limitations}
% =========================================================

The current frontend implementation has the following limitations:

\begin{itemize}
    \item Complaint information is maintained in frontend state and is not
    permanently stored after a page refresh.
    
    \item Worker authentication is currently represented through the
    frontend prototype.
    
    \item Evidence files are selected through the frontend but are not yet
    permanently stored in a backend server.
    
    \item Administrator operations currently focus on viewing and monitoring
    complaint information.
    
    \item Backend API and PostgreSQL database integration are not part of
    this frontend development stage.
\end{itemize}

% =========================================================
\section{Learning Outcomes}
% =========================================================

The development of the Field Worker and Administrator modules provided
practical understanding of:

\begin{itemize}
    \item React component development.
    \item React Router based navigation.
    \item Frontend state management.
    \item Dynamic rendering based on status values.
    \item Interactive progress controls.
    \item Form and input handling.
    \item Basic frontend validation.
    \item File selection and evidence handling.
    \item Dashboard development.
    \item Dynamic calculation of dashboard statistics.
    \item CSS-based interface design.
    \item Frontend testing and debugging.
\end{itemize}


% =========================================================
\end{document}
% =========================================================